\begin{helper}
Let $\mathcal F$ be a finite multi-hypergraph, let $f$ be an edge of
$\mathcal F$, and form the fresh-edge hypergraph
\(\mathcal F_1=\noderef{KUniformFreshEdgeHypergraph}(k,\mathcal F,v)\).
In the line graphs, the independent-triple count in the neighborhood of the
old edge copy does not decrease:
\[
T_{L(\mathcal F)}(f)
\le
T_{L(\mathcal F_1)}(\operatorname{inl} f).
\]
\end{helper}
\begin{proof}
By \noderef{IndependentTripleCount}, the left side counts three-element
subsets of the line-graph neighborhood of $f$ whose three members are
pairwise nonadjacent in \(L(\mathcal F)\).  Define a map from such old
triples to subsets of the new edge-label set by sending every old edge label
$e$ to its old-edge copy \(\operatorname{inl} e\).

This map preserves cardinality three, because the old-edge-copy map is
injective.  It also sends neighbors of $f$ to neighbors of
\(\operatorname{inl} f\): by \noderef{LineGraphOfHypergraph}, old labels
$f$ and $e$ are adjacent exactly when they are distinct and their hyperedges
intersect, and \noderef{KUniformFreshEdgeHypergraph} says old hyperedges in
\(\mathcal F_1\) are precisely the images of the old hyperedges under the
injective vertex map \(x\mapsto \operatorname{inl} x\).  Therefore the
intersection is nonempty before the construction if and only if the
corresponding old-edge-copy intersection is nonempty afterward.

Finally, pairwise nonadjacency inside the triple is preserved for the same
reason: for old labels $e$ and $g$, the copied edges \(\operatorname{inl} e\)
and \(\operatorname{inl} g\) are adjacent in \(L(\mathcal F_1)\) exactly when
$e$ and $g$ are adjacent in \(L(\mathcal F)\).  Thus every old independent
triple maps to a new independent triple.  The map is injective on triples
because \(e\mapsto\operatorname{inl} e\) is injective, so the new counted set
has cardinality at least the old counted set.
\end{proof}
